\section{Details on approximating the predictive variance}
\label{sec:predictiveEntropyAppendix}

We now provide further details on approximating the predictive variance
$v_n(\x|\xopt)$ of inputs $\x$ given the position of the global optimizer
$\xopt$. In particular we include all steps omitted in the presentation of
Section~\ref{sec:predictiveEntropy}.

\subsection{Incorporating the analytic latent constraints (C1.1)}

We first turn to the random variables
\begin{align*}
    \vz &= [f(\xopt);\, \diag[\nabla^2f(\xopt)]], \\
    \vc &= [\vy_n; \nabla f(\xopt); \upper[\nabla^2 f(\xopt)]]
         = [\vy_n; \zero; \zero].
\end{align*}
Here $\vc$ contains the random variables that we will condition on in order to
enforce constraint C1.1. Given the input locations $\x$ and $\xopt$ we can
construct a kernel matrix $\vK$ containing the covariance evaluated on the
stacked vector $[\vz; \vc]$. We again refer to \cite{Solak:2003} in constructing
this matrix which includes derivative observations, the computations of which
are tedious but not overly complicated. Note also that the portions of $\vK$
which correspond to $y_i$ will have an additional $\sigma^2$ due to the
observation noise. Next let $\vK_\vz$, $\vK_\vc$, and $\vK_{\vz\vc}$ denote the
corresponding diagonal and off-diagonal blocks of the kernel matrix.
We can now condition on the observed values of $\vc$ to write
\begin{equation*}
    p(\vz|\D_n,\C{1.1})
    = p(\vz|\vc) = \Normal(\vz|\vm_0,\vV_0)
\end{equation*}
where $\vm_0=\vK_{\vz\vc} \vK^{-1}_\vc\vc$ and $\vV_0= \vK_\vz -
\vK_{\vz\vc}\vK^{-1}_\vc \vK_{\vz\vc}\T$.

\subsection{Incorporating the non-analytic latent constraints (C1.2 and C2)}

The additional constraints C1.2 and C2 can be introduced explicitly as in
(\ref{eq:conditional}), which takes the form of a single Gaussian factor and
$d+1$ non-Gaussian factors
\begin{equation*}
    p(\vz|\D_n,\C1,\C2)
    \propto
    \Normal(\vz|\vm_0,\vV_0)
    \Big[ \prod_{i=1}^{d+1} t_i(z_i) \Big].
\end{equation*}
We approximate this distribution using a single multivariate Gaussian $q(\vz)$
where each non-Gaussian factor is replaced by a Gaussian approximation
$\widetilde t_i(z_i)=\Normal(z_i;\widetilde m_i, \widetilde v_i)$ such that
\begin{equation*}
    q(\vz)
    = \Normal(\vz|\vm,\vV)
    \propto \Normal(\vz|\vm_0,\vV_0)
    \Big[ \prod_{i=1}^{d+1} \Normal(z_i;\widetilde m_i, \widetilde v_i) \Big]
\end{equation*}
where this approximation is parameterized by
$\vm=\vV[\widetilde\vV^{-1}\widetilde\vm+\vV_0^{-1}\vm_0]$ and
$\vV=(\widetilde\vV^{-1}+\vV_0^{-1})^{-1}$. The parameters of the approximate
factors are combined to form the vector $[\widetilde\vm]_i=\widetilde m_i$ and
the diagonal matrix $[\widetilde\vV]_{ii}=\widetilde v_i$. 

\def\oldm{\bar m_{i}}
\def\oldv{\bar v_{i}}
\def\norm{\bar Z_{i}}

To compute the approximate factors we use expectation propagation (EP). EP is a
procedure that starts from some initial values for the approximate factors
$(\widetilde m_i, \widetilde v_i)$ and iteratively refines these quantities;
here we initialize $\widetilde m_i=0$ and $\widetilde v_i=\infty$ which
corresponds to $\vm=\vm_0$ and $\vV=\vV_0$.  At each iteration, for every factor
$i$, we remove the contribution of the $i$th approximate factor to form the
\emph{cavity} distribution $q_{\setminus i}(\vz)\propto q(\vz)/\widetilde
t_i(z_i)$. Given the independent factors we consider here we can focus on each
individual component $q_{\setminus i}(z_i)$ separately with mean and variance
\begin{align*}
    \oldm &= \oldv(m_i/v_{ii} - \widetilde m_i/\widetilde v_i), \\
    \oldv &= (v_{ii}^{-1} - \widetilde v_i^{-1})^{-1}.
\end{align*}
Let $\hat q(z_i)\propto q_{\setminus i}(z_i) t_i(z_i)$ denote the \emph{tilted}
distribution where the $i$th approximate factor has been replaced by the
corresponding real factor. EP proceeds by finding the approximation $q_i$ that
minimizes the KL-divergence $\mathrm{D}[\hat q_i||q_i]$ where $q_i$ is
restricted to be Gaussian. This amounts to matching the first two moments.
Finally, by removing the influence of the cavity distribution and setting
$\widetilde t_i(z_i)\propto q_i(z_i)/q_{\setminus i}(z_i)$ we can update the
approximate factors. This can be performed using the same procedure which forms
the cavity distribution.

For both sets of constraints used in this work the moments can easily be
obtained by computing the normalizing constant $\norm = \int
\Normal(z_i|\oldm,\oldv) \,t_i(z_i) \,dz_i$ and using the following identities:
\begin{align*}
    \E_{\hat q}(z_i)
    &= \oldm + \oldv \frac{\partial\norm}{\partial\oldm},
    &
    \textrm{Var}_{\hat q}(z_i)
    &= \oldv - \oldv^2
    \left(
    \frac{\partial\norm}{\partial\oldm} 
    - 2\frac{\partial\norm}{\partial\oldv}\right).
\end{align*}
For the factors corresponding to constraints on the diagonal Hessian, i.e.\
where $t_i(z_i)=\I[z_i<0]$, the distribution is also simply a truncated
Gaussian.  Given these moments we can remove the contribution of the cavity
distribution as above and write
\begin{align*}
    \widetilde m_i &\gets \oldm - \kappa^{-1},
    &\text{where } 
      \alpha &= -\frac{\oldm}{\sqrt{\oldv}}, \\
    \widetilde v_i &\gets \beta^{-1} - \oldv,
    & \beta &=
      \frac{\phi(\alpha)}{\Phi(\alpha)}\left[\frac{\phi(\alpha)}{\Phi(\alpha)} +
      \alpha\right] \frac1{\oldv}, \\
    &&
    \kappa &= \left[\frac{\phi(\alpha)}{\Phi(\alpha)}-\alpha\right]\frac1{\sqrt{\oldv}}.
\end{align*}
For the final soft-maximum constraint,
$\Phi\big((z_i-y_\mathrm{max})/\sigma\big)$, the moments can be calculated in a
similar fashion. Using the same procedure as above we arrive at very similar
updates:
\begin{align*}
    \widetilde m_i &\gets \oldm + \kappa^{-1},
    &\text{where } 
      \alpha &= \frac{\oldm-y_\mathrm{max}}{\sqrt{\oldv+\sigma^2}}, \\
    \widetilde v_i &\gets \beta^{-1} - \oldv,
    & \beta &=
      \frac{\phi(\alpha)}{\Phi(\alpha)}\left[\frac{\phi(\alpha)}{\Phi(\alpha)} +
      \alpha\right] \frac1{\oldv+\sigma^2}, \\
    &&
    \kappa &= \left[\frac{\phi(\alpha)}{\Phi(\alpha)}+\alpha\right]\frac1{\sqrt{\oldv+\sigma^2}}.
\end{align*}

\subsection{Incorporating the prediction constraint (C3)}

Given some test input $\x$ we now turn to the problem of making predictions
about $f(\x)$. We again note that both the ``prior'' terms $\vm_0$, $\vV_0$ and
the EP factors, $\widetilde\vm$ and $\widetilde\vV$, are independent
of $\x$ and can be precomputed once for later use at prediction time. 

Let $\vf=[f(\vx); f(\xopt)]$ be a vector given by the concatenation of the
latent function at $\x$ and $\xopt$. The distribution for $\vf$ given the first
two constraints can be written as 
\begin{equation}
    p(\vf|\D_n,\C1,\C2)
    \approx
    \int p(\vf|\vz,\vc) \,q(\vz) \,d\vz = \Normal(\vf|\vm_\vf,\vV_\vf)\,.
\end{equation}
By writing $p(\vf|\vz,\vc)$ above we are assuming that $\vf$ is independent of
C1.2 and C2 given $\vz$ and as a result the above is simply an integral over the
product of two Gaussians. Let $\vK_\dagger$ be the cross-covariance matrix
evaluated between $\vf$ and $[\vz;\vc]$ and $\vK_\vf$ the covariance matrix
associated with $\vf$.  The posterior above will then be Gaussian with mean and
variance
\begin{align*}
    \vm_\vf &= \vK_\dagger [ \vK + \widetilde\vW]^{-1} [\vc; \widetilde\vm] \\
    \vV_\vf &= \vK_\vf - \vK_\dagger [\vK + \widetilde\vW]^{-1}\vK_\dagger\T,
\end{align*}
where $\widetilde\vW$ is a block-diagonal matrix where the first block is zero
and the second is $\widetilde\vV$ (note this matrix is also diagonal since
$\widetilde\vV$ is diagonal). Finally, these values can be plugged into
(\ref{eq:finalPredictive}--\ref{eq:predictiveVariance}) in order to arrive at
$v_n(\x|\xopt)$.

